Orbiter can classify cubic curves in $\PG(2,q)$. 
To this end, the $(9,3)$-arcs in $\PG(2,q)$ are classified first.
From this classification, the classification of curves is computed.
This classification only works for arcs which contain a $(9,3)$ arc. 
For very small fields, this is not always the case. 
Here is an example. The command  

\sourcecode{cubic_curves_PG_2_4}

classifies the cubic curves in $\PG(2,4)$ containing at least 9 points.
There are none.
%The command 

%\sourcecode{cubic_curves_PG_2_4_draw}

%draws the poset of orbits of the classification of 9-arcs, shown below:

%\orbiteroutput{GRAPHICS/WRAPPERS/cubic_curve_arcs_q4}


\bigskip


The command  

\sourcecode{poly_orbits_q4}

uses the basic Schreier algorithm to classify all cubic curves in $\PG(2,4).$
This is done using the action of the group $\PGGL(3,4)$ on all cubic forms over $\bbF_4$.
The output is shown below:

\orbiteroutput{GRAPHICS/LATEX/poly_orbits_d3_n2_q4}

This classification involves curves that are degenerate or singular, or contain lines.